\chapter{Supersymmetric Gauge Theory}
\label{ch:susy-gauge-theory}

\ConventionsLorentzian

This chapter introduces supersymmetric gauge theory as field-variable data.
The particle spectrum, vacuum structure, Wilsonian effective action, and
nonperturbative dynamics are later constructions.  We keep the gauge
representative, the gauge-invariant observables, and the regularization
scheme separate.

\section{Gauge and Trace Conventions}
\label{sec:susy-gauge-trace-conventions}

This chapter uses the Hermitian-generator convention of
Chapter~\ref{ch:vol2-classical-yang-mills-matter}.  If
\(\operatorname{Lie}(G)\) is realized by anti-Hermitian matrices in a unitary
representation, write
\[
  \mathfrak g:=\ii\,\operatorname{Lie}(G)
\]
for the corresponding real vector space of Hermitian generators.  Its compact
Lie bracket in Hermitian coordinates is
\[
  [X,Y]_{\rm H}:=-\ii[X,Y],
  \qquad X,Y\in\mathfrak g,
\]
where the bracket on the right is the ordinary matrix commutator.  Fix an
invariant nondegenerate symmetric bilinear form
\(\operatorname{tr}:\mathfrak g\otimes\mathfrak g\to\R\).  In trace-delta
examples one chooses a basis \(t^a\) with
\(\operatorname{tr}(t^a t^b)=\delta^{ab}\).

In a local trivialization, a gauge field is a Hermitian
\(\mathfrak g\)-valued one-form \(A=A_\mu dx^\mu\), and the covariant
derivative on a field in a unitary representation is
\[
  \nabla_\mu=\partial_\mu-\ii A_\mu .
\]
The curvature is
\[
  F_{\mu\nu}
  :=
  \ii[\nabla_\mu,\nabla_\nu]
  =
  \partial_\mu A_\nu-\partial_\nu A_\mu-\ii[A_\mu,A_\nu].
\]
Under a smooth gauge transformation \(u:U\to G\),
\[
  A_\mu^u=uA_\mu u^{-1}+\ii u\,\partial_\mu u^{-1}.
\]

\begin{proposition}[Curvature covariance and gauge-invariant densities]
\label{prop:susy-gauge-curvature-covariance}
With the definitions above,
\[
  F_{\mu\nu}^u=uF_{\mu\nu}u^{-1}.
\]
Consequently
\(\operatorname{tr}(F_{\mu\nu}F^{\mu\nu})\) and
\(\operatorname{tr}(F_{\mu\nu}\widetilde F^{\mu\nu})\) are local
gauge-invariant scalar densities.
\end{proposition}

\begin{proof}
The transformation law for \(A_\mu\) is equivalent to
\[
  \nabla_\mu^u=u\nabla_\mu u^{-1}
\]
as an equality of differential operators.  Therefore
\[
  F_{\mu\nu}^u
  =
  \ii[\nabla_\mu^u,\nabla_\nu^u]
  =
  \ii u[\nabla_\mu,\nabla_\nu]u^{-1}
  =
  uF_{\mu\nu}u^{-1}.
\]
Invariance of the bilinear form gives
\(\operatorname{tr}(uXu^{-1}uYu^{-1})=\operatorname{tr}(XY)\).  Applying
this to \(X=F_{\mu\nu}\) and \(Y=F^{\mu\nu}\), and then to
\(Y=\widetilde F^{\mu\nu}\), proves the two density statements.
\end{proof}

The Yang--Mills coupling convention in this monograph is
\[
  \mathcal L_{\rm YM}
  =
  -\frac{1}{4g^2}\operatorname{tr}(F_{\mu\nu}F^{\mu\nu})
  +\frac{\theta}{32\pi^2}
  \operatorname{tr}(F_{\mu\nu}\widetilde F^{\mu\nu})
  +\cdots ,
\]
where \(\widetilde F^{\mu\nu}=\frac12
\epsilon^{\mu\nu\rho\sigma}F_{\rho\sigma}\) in the stated orientation
convention.  This convention applies uniformly to \(U(1)\), \(U(N)\), and
simple factors after the bilinear form has been fixed.

\section{Vector Multiplet as Field Variables}

\begin{definition}[Off-shell vector multiplet]
\label{def:off-shell-vector-multiplet}
An \(\mathcal N=1\) off-shell vector multiplet with gauge algebra
\(\mathfrak g\) consists locally of a connection \(A_\mu\), a Weyl spinor
\(\lambda_\alpha\) valued in \(\mathfrak g_\C\), its conjugate
\(\bar\lambda_{\dot\alpha}\), and a real auxiliary field \(D\) valued in
\(\mathfrak g\).  These are gauge-representative field variables.  The
physical local observables are gauge-invariant cohomological or operator
classes after quantization.
\end{definition}

In Wess--Zumino gauge, the supersymmetry transformations have the form
\[
  \delta A_\mu
  =
  \ii\epsilon\sigma_\mu\bar\lambda
  -\ii\lambda\sigma_\mu\bar\epsilon ,
\]
\[
  \delta\lambda_\alpha
  =
  (\sigma^{\mu\nu}\epsilon)_\alpha F_{\mu\nu}
  +\ii\epsilon_\alpha D,
  \qquad
  \delta D
  =
  -\epsilon\sigma^\mu D_\mu\bar\lambda
  -D_\mu\lambda\sigma^\mu\bar\epsilon .
\]
Here \(D_\mu X:=\partial_\mu X-\ii[A_\mu,X]\) for adjoint-valued component
fields \(X\).  The displayed transformations are not bare supertranslations
restricted to the Wess--Zumino slice.  They are bare supertranslations
followed by the compensating chiral gauge transformation that restores the
representative \eqref{eq:wess-zumino-gauge-vector-superfield}.

\begin{proposition}[Closure in Wess--Zumino gauge]
\label{prop:wess-zumino-gauge-closure}
Let \(\mathcal V\) be a local formal neighbourhood in the space of real
vector superfields, and let \(\mathcal S\subset\mathcal V\) be the
Wess--Zumino slice of
Proposition~\ref{prop:wess-zumino-gauge-representative}.  Assume the local
slice and the compensating chiral gauge parameters are chosen in this
neighbourhood.  For each supersymmetry parameter \(\epsilon\), write
\[
  \widehat\delta_\epsilon
  =
  \delta^0_\epsilon+\delta^{\rm g}_{\Lambda_\epsilon}
\]
for the vector field tangent to \(\mathcal S\), where
\(\delta^0_\epsilon\) is the supertranslation vector field on
\(\mathcal V\), and \(\delta^{\rm g}_{\Lambda_\epsilon}\) is the chiral
gauge vector field with parameter \(\Lambda_\epsilon(V)\).  Then, on
\(\mathcal S\),
\[
  [\widehat\delta_{\epsilon_1},\widehat\delta_{\epsilon_2}]
  =
  a^\rho\partial_\rho+\delta^{\rm g}_{\Omega_{12}},
  \qquad
  a^\rho=
  2\ii\left(
  \epsilon_1\sigma^\rho\bar\epsilon_2
  -\epsilon_2\sigma^\rho\bar\epsilon_1
  \right),
  \label{eq:wess-zumino-gauge-closure}
\]
where \(\Omega_{12}\) is a real residual ordinary gauge parameter determined
by the two compensators and the slice condition.  Equivalently, if
\(\Xi_{12}:=\Omega_{12}+a^\rho A_\rho\), then for every adjoint-valued
component field \(X\in\{\lambda,\bar\lambda,D\}\),
\[
  [\widehat\delta_{\epsilon_1},\widehat\delta_{\epsilon_2}]X
  =
  a^\rho D_\rho X+\ii[\Xi_{12},X],
\]
and for the connection
\[
  [\widehat\delta_{\epsilon_1},\widehat\delta_{\epsilon_2}]A_\mu
  =
  a^\rho F_{\rho\mu}+D_\mu\Xi_{12}.
\]
No Euler--Lagrange equation is used; the closure is off shell because the
auxiliary field \(D\) is part of the coordinate system.
\end{proposition}

\begin{proof}
On the full real-vector-superfield space \(\mathcal V\), the
supertranslation vector fields obey the algebra derived in
Chapter~\ref{ch:susy-superspace-superfields-local-actions}:
\[
  [\delta^0_{\epsilon_1},\delta^0_{\epsilon_2}]
  =
  a^\rho\partial_\rho .
\]
Let \(\mathcal D_{\rm g}\) be the vertical distribution generated by chiral
gauge transformations.  This distribution is involutive because chiral gauge
transformations form a group, and it is preserved by supertranslations:
if \(\bar D_{\dot\alpha}\Lambda=0\), then the supertranslation of
\(\Lambda\) is again a chiral gauge parameter.  Field dependence of
\(\Lambda_\epsilon(V)\) does not change this conclusion; differentiating a
chiral gauge parameter along any tangent vector in field space still gives a
chiral gauge parameter.  Therefore every term in the commutator
\[
  [\delta^0_{\epsilon_1}+\delta^{\rm g}_{\Lambda_{\epsilon_1}},
    \delta^0_{\epsilon_2}+\delta^{\rm g}_{\Lambda_{\epsilon_2}}]
\]
except the bare supertranslation commutator is vertical.  Hence the
commutator of the projected vector fields equals \(a^\rho\partial_\rho\)
modulo a chiral gauge vector field.

Both \(\widehat\delta_{\epsilon_i}\) are tangent to the Wess--Zumino slice,
so their commutator is tangent to that slice.  A vertical vector field that
is tangent to the Wess--Zumino slice is precisely the residual ordinary
gauge transformation left after the components \(C,\chi,M\) of a general
real vector superfield have been removed.  This gives
\eqref{eq:wess-zumino-gauge-closure}.  In Hermitian-generator conventions an
infinitesimal residual gauge parameter \(\Omega\) acts by
\[
  \delta^{\rm g}_\Omega A_\mu=D_\mu\Omega,
  \qquad
  \delta^{\rm g}_\Omega X=\ii[\Omega,X]
\]
on adjoint component fields.  The alternative covariant form follows from
the identity
\[
  a^\rho\partial_\rho A_\mu+D_\mu\Omega
  =
  a^\rho F_{\rho\mu}+D_\mu(\Omega+a^\rho A_\rho),
\]
which uses
\(F_{\rho\mu}=\partial_\rho A_\mu-\partial_\mu A_\rho
-\ii[A_\rho,A_\mu]\) and
\(D_\mu Y=\partial_\mu Y-\ii[A_\mu,Y]\).  This proves the stated closure
without a component equation of motion.
\end{proof}

Thus Wess--Zumino gauge is not a gauge-invariant description of
supersymmetry; it is a convenient representative after ordinary gauge
freedom has been partly fixed.

\section{Superfield Strength}

Let \(V\) be a real vector superfield.  For matrix notation in a
representation of \(G\), a chiral gauge transformation acts by
\[
  e^V\longmapsto e^{-\bar\Lambda}e^V e^{-\Lambda},
  \qquad \bar D_{\dot\alpha}\Lambda=0 .
\]
The chiral field strength is
\[
  W_\alpha
  =
  -\frac14\bar D^2\!\left(e^{-V}D_\alpha e^V\right),
  \qquad
  \bar D_{\dot\alpha}W_\beta=0.
\]
It transforms covariantly,
\[
  W_\alpha\longmapsto e^\Lambda W_\alpha e^{-\Lambda}.
\]
Consequently \(\operatorname{tr}(W^\alpha W_\alpha)\) is a chiral
gauge-invariant local superfield.  With the trace convention fixed above, the
gauge kinetic term is normalized so that its component expansion contains the
Yang--Mills term
\[
  -\frac{1}{4g^2}\operatorname{tr}(F_{\mu\nu}F^{\mu\nu}).
\]
All later anomaly and beta-function formulae must refer back to this coupling
normalization.

\section{Matter Multiplets and the Moment Map}

Let \(R\) be a finite-dimensional unitary complex representation of \(G\).
A chiral matter multiplet is a chiral superfield \(\Phi\) valued in \(R\).
Gauge covariance of the kinetic term is expressed by
\[
  \int d^4\theta\, \bar\Phi e^V\Phi .
\]
In components this term contains the scalar kinetic term, fermion kinetic
term, gauge couplings, and the auxiliary coupling
\[
  \langle\mu_R(\phi,\bar\phi),D\rangle,
\]
where \(\mu_R\in\mathfrak g^*\) is the moment map determined by
\[
  \langle \mu_R(\phi,\bar\phi),X\rangle
  =
  \phi^\dagger X_R\phi,
  \qquad X\in\mathfrak g .
\]
Here \(X_R\) is the Hermitian generator acting in \(R\).  A
Fayet--Iliopoulos parameter is a linear functional
\(\zeta\in\mathfrak g^*\) satisfying
\[
  \zeta([X,Y]_{\rm H})=0,\qquad X,Y\in\mathfrak g,
\]
so it factors through the abelian quotient of \(\mathfrak g\).  In
particular there is no FI parameter for a semisimple gauge algebra.

\begin{proposition}[Auxiliary \(D\)-field elimination]
\label{prop:susy-d-term-auxiliary-elimination}
Let \(\eta=\mu_R(\phi,\bar\phi)-\zeta\), with \(\zeta=0\) if no FI parameter
is allowed.  Use the bilinear form to define
\(\eta^\sharp\in\mathfrak g\) by
\[
  \operatorname{tr}(\eta^\sharp X)=\langle\eta,X\rangle
  \quad\text{for every }X\in\mathfrak g .
\]
For fixed scalar fields, the auxiliary part of the Lorentzian Lagrangian is
\[
  \mathcal L_D
  =
  \frac{1}{2g^2}\operatorname{tr}(D^2)
  +\langle\eta,D\rangle .
\]
Eliminating \(D\) gives
\[
  D=-g^2\eta^\sharp,
  \qquad
  \mathcal L_D\big|_{\rm on\ shell}
  =
  -\frac{g^2}{2}\operatorname{tr}(\eta^\sharp\eta^\sharp).
\]
Thus the corresponding scalar potential contribution is
\[
  V_D=\frac{g^2}{2}\operatorname{tr}(\eta^\sharp\eta^\sharp)
\]
in this convention.
\end{proposition}

\begin{proof}
By the definition of \(\eta^\sharp\),
\[
  \mathcal L_D
  =
  \frac{1}{2g^2}\operatorname{tr}(D^2)
  +\operatorname{tr}(\eta^\sharp D).
\]
Completing the square gives
\[
  \mathcal L_D
  =
  \frac{1}{2g^2}
  \operatorname{tr}\!\left(D+g^2\eta^\sharp\right)^2
  -\frac{g^2}{2}\operatorname{tr}(\eta^\sharp\eta^\sharp).
\]
The algebraic Euler--Lagrange equation is therefore
\(D+g^2\eta^\sharp=0\).  Substitution gives the displayed on-shell
Lagrangian.  Since Lorentzian scalar Lagrangians have the form kinetic minus
potential, the potential is the positive final expression.
\end{proof}

\section{Gauge Anomalies as a Consistency Condition}

The classical field-variable construction is not yet a quantum gauge theory.
For a four-dimensional chiral gauge theory, gauge invariance of the quantum
theory requires the perturbative gauge anomaly class to vanish.  In a
representation \(R\), the cubic anomaly coefficient is the symmetric
trilinear form
\[
  A_R(X,Y,Z)
  =
  \operatorname{Tr}_R\!\left(X_R\{Y_R,Z_R\}\right),
  \qquad X,Y,Z\in\mathfrak g ,
\]
with the trace over chiral fermion variables.  The local anomaly cancellation
condition is
\[
  \sum_{\text{chiral fermions }i} A_{R_i}=0
\]
as an invariant symmetric trilinear form.  Global anomalies and mixed
anomalies require separate topological data and are not detected by this
polynomial alone.

\begin{proposition}[Conjugate-pair and real-representation cancellation]
\label{prop:susy-gauge-anomaly-conjugate-cancellation}
Let \(R^\vee\) be the conjugate representation of a unitary representation
\(R\), in the Hermitian-generator convention.  Then
\[
  A_{R^\vee}=-A_R .
\]
Consequently a vectorlike pair \(R\oplus R^\vee\) has zero perturbative local
gauge anomaly.  If \(R\) is equivalent to its conjugate representation, then
\(A_R=0\).  In particular, the adjoint gaugino in an \(\mathcal N=1\) vector
multiplet gives no perturbative cubic gauge anomaly.
\end{proposition}

\begin{proof}
If \(T_R^a\) are Hermitian generators in \(R\), then the conjugate
representation has Hermitian generators
\[
  T_{R^\vee}^a=-(T_R^a)^* .
\]
Therefore
\[
  \operatorname{Tr}_{R^\vee}
  \left(T_{R^\vee}^a\{T_{R^\vee}^b,T_{R^\vee}^c\}\right)
  =
  -
  \left[
  \operatorname{Tr}_{R}
  \left(T_R^a\{T_R^b,T_R^c\}\right)
  \right]^* .
\]
The trace on the right is real because the anticommutator of Hermitian
matrices is Hermitian and the invariant symmetric tensor is obtained by
symmetrizing over \(a,b,c\).  Hence \(A_{R^\vee}=-A_R\).  The vectorlike
statement follows by addition.  If \(R\simeq R^\vee\), anomaly coefficients
are representation invariants, so \(A_R=A_{R^\vee}=-A_R\), and \(A_R=0\).
The adjoint representation of a compact Lie algebra is real, giving the last
claim.
\end{proof}

\begin{openproblem}[Supersymmetric gauge theory with regulated Ward identities]
\label{op:susy-gauge-regulated-ward-identities}
Construct four-dimensional supersymmetric gauge theories with a regulator
that states the gauge symmetry, supersymmetry transformations, BV master
equation, local counterterm space, anomaly class, and decoupling limit at
finite cutoff.  This is the regulator-level input needed before exact
Wilsonian statements about holomorphy, moduli spaces, and duality can be
treated as quantum statements rather than coordinate calculations.
\end{openproblem}
